\documentclass[11pt,a4paper]{article}
\usepackage[utf8]{inputenc}
\usepackage[T1]{fontenc}
\usepackage[spanish,es-minimal,es-noshorthands]{babel}
\usepackage{lmodern}
\usepackage[sfdefault,lining]{FiraSans}
\usepackage[a4paper,top=2.3cm,bottom=2.2cm,left=1.8cm,right=1.8cm,headheight=1.3cm,headsep=0.4cm,footskip=1cm]{geometry}
\usepackage{amsmath,amssymb}
\usepackage{icomma}
\usepackage{graphicx}
\usepackage[table]{xcolor}
\usepackage{tikz}
\usepackage{fontawesome5}
\usepackage{ccicons}
\usepackage{fancyhdr}
\usepackage{lastpage}
\usepackage{tasks}
\usepackage{needspace}
\usepackage[most]{tcolorbox}
\usepackage{array}
\usepackage[hidelinks]{hyperref}
\AtBeginDocument{\hypersetup{pdftitle={Ficha de ejercicios de la unidad 3: Proporcionalidad y porcentajes. 3º ESO},pdfsubject={Matemáticas 3º ESO},pdfkeywords={porcentajes, aumentos, disminuciones, interés compuesto, repartos, mezclas, movimientos, proporcionalidad compuesta}}}

\definecolor{navy}{HTML}{1B2748}
\definecolor{accent}{HTML}{E8603A}
\definecolor{teal}{HTML}{1F9E90}
\definecolor{soft}{HTML}{EEF1F8}
\definecolor{muted}{HTML}{5D6685}

\setlength{\parindent}{0pt}
\setlength{\parskip}{3pt}
\settasks{label=\alph*), label-format=\bfseries\color{navy}, label-width=1.3em, item-indent=1.8em, column-sep=0.6em, after-item-skip=2pt, before-skip=3pt, after-skip=4pt}

% ---------- capçalera i peu ----------
\newcommand{\llicencia}{\href{https://creativecommons.org/licenses/by-nc-sa/4.0/deed.es}{\ccbyncsa}\ \ Licencia Creative Commons Atribución-NoComercial-CompartirIgual 4.0 Internacional (CC BY-NC-SA 4.0)}
\pagestyle{fancy}
\fancyhf{}
\fancyhead[L]{\raisebox{-0.35\height}{\includegraphics[height=0.95cm]{logo-gva.png}}}
\fancyhead[C]{\small\color{muted}Ficha de ejercicios · Unidad 3: Proporcionalidad y porcentajes · 3\textsuperscript{o} ESO}
\fancyhead[R]{\raisebox{-0.35\height}{\includegraphics[height=0.95cm]{logo-vermellar.png}}}
\fancyfoot[L]{\scriptsize\color{muted}\llicencia}
\fancyfoot[R]{\small\color{muted}\thepage}
\renewcommand{\headrulewidth}{0.4pt}
\renewcommand{\headrule}{\color{muted!40}\hrule height \headrulewidth}
\fancypagestyle{primera}{\fancyhf{}\renewcommand{\headrulewidth}{0pt}
  \fancyfoot[L]{\scriptsize\color{muted}\llicencia}
  \fancyfoot[R]{\small\color{muted}\thepage}}

% ---------- nivells i exercicis ----------
\newcommand{\estrelles}[1]{{\color{accent}\ifcase#1\or\faStar\faStar[regular]\faStar[regular]\or\faStar\faStar\faStar[regular]\or\faStar\faStar\faStar\fi}}
\newcounter{exer}
\newcommand{\ex}[1]{\par\needspace{4\baselineskip}\medskip\stepcounter{exer}\noindent{\color{navy}\bfseries\theexer.}\ \ {\small\estrelles{#1}}\ \ }

\newcommand{\seccio}[2]{\par\needspace{10\baselineskip}\bigskip
  \begin{tcolorbox}[colback=navy,colframe=navy,coltext=white,arc=2mm,boxrule=0pt,left=3mm,right=3mm,top=1.2mm,bottom=1.2mm]
  {\large\bfseries #1\quad #2}
  \end{tcolorbox}\vspace{-2pt}}

% ---------- recta real ----------
\newcommand{\rectabuida}{\begin{tikzpicture}[x=0.42cm,baseline=0]
  \draw[-latex,muted] (-6.7,0)--(6.8,0);
  \foreach \i in {-6,...,6}{\draw[muted] (\i,0.12)--(\i,-0.12) node[below=-1pt,font=\scriptsize,text=muted]{$\i$};}
  \end{tikzpicture}}
% \rectaint{x1}{x2}{estil esquerra}{estil dreta}: estil = ple, buit, inf
\newcommand{\extrem}[2]{\ifx#2p\fill[accent] (#1,0) circle (0.13);\else\ifx#2b\filldraw[fill=white,draw=accent,line width=1.2pt] (#1,0) circle (0.13);\fi\fi}
\newcommand{\rectaint}[4]{\begin{tikzpicture}[x=0.42cm,baseline=0]
  \draw[-latex,muted] (-6.7,0)--(6.8,0);
  \foreach \i in {-6,...,6}{\draw[muted] (\i,0.12)--(\i,-0.12) node[below=-1pt,font=\scriptsize,text=muted]{$\i$};}
  \draw[accent,line width=3pt] (#1,0)--(#2,0);
  \ifx#3i\fill[accent] (#1-0.25,0)--(#1+0.15,0.2)--(#1+0.15,-0.2)--cycle;\else\extrem{#1}{#3}\fi
  \ifx#4i\fill[accent] (#2+0.25,0)--(#2-0.15,0.2)--(#2-0.15,-0.2)--cycle;\else\extrem{#2}{#4}\fi
  \end{tikzpicture}}

\newcolumntype{C}[1]{>{\centering\arraybackslash}p{#1}}

% =====================================================================
% Capçalera FentMates de les fitxes: un color per a cada curs.
% Els colors són els de la pàgina inicial (LEVELS a assets/js/materials.js);
% el logo ve de _eines/logo/pdf/ (python3 _eines/logo/genera.py).
%
% Ús, al final del preàmbul de la unitat:
%   % =====================================================================
% Capçalera FentMates de les fitxes: un color per a cada curs.
% Els colors són els de la pàgina inicial (LEVELS a assets/js/materials.js);
% el logo ve de _eines/logo/pdf/ (python3 _eines/logo/genera.py).
%
% Ús, al final del preàmbul de la unitat:
%   % =====================================================================
% Capçalera FentMates de les fitxes: un color per a cada curs.
% Els colors són els de la pàgina inicial (LEVELS a assets/js/materials.js);
% el logo ve de _eines/logo/pdf/ (python3 _eines/logo/genera.py).
%
% Ús, al final del preàmbul de la unitat:
%   \input{../comu/fentmates.tex}
%   \fmcurs{1eso}
% i, al començament del document:
%   \capcalera{Matemàtiques}{1r ESO}{Fitxa d'exercicis · Unitat 2}{Divisibilitat}
% Opcional abans de \input: \def\fmneed{14} (espai mínim abans de cada secció).
% =====================================================================
\makeatletter
\def\fm@logodir{../../logo/pdf}
\providecommand{\fmneed}{10}
\def\fm@cursnom{}

% clau del curs, color i color del text damunt del color
\newcommand{\fm@curs}[3]{\@namedef{fm@color@#1}{#2}\@namedef{fm@text@#1}{#3}}
\fm@curs{1eso}{E6B84A}{navy}
\fm@curs{2eso}{3B82F6}{white}
\fm@curs{3eso}{FF7A50}{white}
\fm@curs{4eso-a}{A78BFA}{white}
\fm@curs{4eso-b}{6D28D9}{white}
\fm@curs{1batx-mi}{9B1C1C}{white}
\fm@curs{1batx-ccss}{EF6B6B}{white}
\fm@curs{2batx-mii}{1E6B47}{white}
\fm@curs{2batx-ccss}{2BB3A3}{white}

\newcommand{\fmcurs}[1]{%
  \@ifundefined{fm@color@#1}{\PackageError{fentmates}{Curs desconegut: #1}{}}{}%
  \edef\fm@tmp{\noexpand\definecolor{curs}{HTML}{\@nameuse{fm@color@#1}}%
               \noexpand\colorlet{curstext}{\@nameuse{fm@text@#1}}}\fm@tmp
  \colorlet{cursfosc}{curs!72!black}%
  \def\fm@logo{\fm@logodir/vertical-#1.pdf}%
  \def\fm@marca{\fm@logodir/marca-#1.pdf}}

% ---------- capçalera de la primera pàgina ----------
% #1 matèria (petit), #2 curs, #3 tipus de document i unitat, #4 títol
\newcommand{\capcalera}[4]{%
  \gdef\fm@cursnom{#2}%
  \par\noindent
  \begin{tcolorbox}[colback=curs,colframe=curs,boxrule=0pt,arc=2.2mm,width=3.5cm,height=2.15cm,
      valign=center,left=3mm,right=2mm,top=1mm,bottom=1mm,nobeforeafter,box align=center]
    \color{curstext}{\footnotesize #1}\par\vspace{1pt}{\fontsize{19}{21}\selectfont\bfseries #2\par}%
  \end{tcolorbox}\hspace{5mm}%
  \begin{minipage}[c]{\dimexpr\linewidth-3.5cm-5mm-2.9cm\relax}
    {\color{cursfosc}\small\bfseries\MakeUppercase{#3}}\par\vspace{2pt}
    {\color{navy}\fontsize{20}{23}\selectfont\bfseries #4\par}
  \end{minipage}\hfill
  \begin{minipage}[c]{2.6cm}\raggedleft\includegraphics[width=2.6cm]{\fm@logo}\end{minipage}\par
  \vspace{2pt}}

% ---------- barres de secció del color del curs ----------
\renewcommand{\seccio}[2]{\par\needspace{\fmneed\baselineskip}\bigskip
  \begin{tcolorbox}[colback=curs,colframe=curs,coltext=curstext,arc=2mm,boxrule=0pt,left=3mm,right=3mm,top=1.2mm,bottom=1.2mm]
  {\large\bfseries #1\quad #2}
  \end{tcolorbox}\vspace{-2pt}}

% ---------- estrelles de dificultat del color del curs (a Batxillerat no se'n posen) ----------
\@ifundefined{estrelles}{}{%
  \renewcommand{\estrelles}[1]{{\color{curs}\ifcase#1\or\faStar\faStar[regular]\faStar[regular]\or\faStar\faStar\faStar[regular]\or\faStar\faStar\faStar\fi}}}

% ---------- capçalera i peu de les pàgines ----------
% primera pàgina: logos institucionals al peu
\fancypagestyle{primera}{\fancyhf{}\renewcommand{\headrulewidth}{0pt}
  \fancyfoot[L]{\raisebox{-0.32\height}{\includegraphics[height=0.78cm]{logo-gva.png}}\hspace{5mm}%
    \raisebox{-0.32\height}{\includegraphics[height=0.78cm]{logo-vermellar.png}}}
  \fancyfoot[R]{\scriptsize\color{muted}\href{https://creativecommons.org/licenses/by-nc-sa/4.0/deed.ca}{\ccbyncsa}\ \ CC BY-NC-SA 4.0 · fentmates.org\qquad\small\thepage}}
% la resta: el símbol a l'esquerra i el curs a la dreta, amb el títol al centre com abans
\fancyhead[L]{\raisebox{-0.3\height}{\includegraphics[height=0.62cm]{\fm@marca}}}
\fancyhead[R]{\small\bfseries\color{cursfosc}\fm@cursnom}
\renewcommand{\headrule}{\color{curs!45}\hrule height \headrulewidth}
\makeatother

%   \fmcurs{1eso}
% i, al començament del document:
%   \capcalera{Matemàtiques}{1r ESO}{Fitxa d'exercicis · Unitat 2}{Divisibilitat}
% Opcional abans de \input: \def\fmneed{14} (espai mínim abans de cada secció).
% =====================================================================
\makeatletter
\def\fm@logodir{../../logo/pdf}
\providecommand{\fmneed}{10}
\def\fm@cursnom{}

% clau del curs, color i color del text damunt del color
\newcommand{\fm@curs}[3]{\@namedef{fm@color@#1}{#2}\@namedef{fm@text@#1}{#3}}
\fm@curs{1eso}{E6B84A}{navy}
\fm@curs{2eso}{3B82F6}{white}
\fm@curs{3eso}{FF7A50}{white}
\fm@curs{4eso-a}{A78BFA}{white}
\fm@curs{4eso-b}{6D28D9}{white}
\fm@curs{1batx-mi}{9B1C1C}{white}
\fm@curs{1batx-ccss}{EF6B6B}{white}
\fm@curs{2batx-mii}{1E6B47}{white}
\fm@curs{2batx-ccss}{2BB3A3}{white}

\newcommand{\fmcurs}[1]{%
  \@ifundefined{fm@color@#1}{\PackageError{fentmates}{Curs desconegut: #1}{}}{}%
  \edef\fm@tmp{\noexpand\definecolor{curs}{HTML}{\@nameuse{fm@color@#1}}%
               \noexpand\colorlet{curstext}{\@nameuse{fm@text@#1}}}\fm@tmp
  \colorlet{cursfosc}{curs!72!black}%
  \def\fm@logo{\fm@logodir/vertical-#1.pdf}%
  \def\fm@marca{\fm@logodir/marca-#1.pdf}}

% ---------- capçalera de la primera pàgina ----------
% #1 matèria (petit), #2 curs, #3 tipus de document i unitat, #4 títol
\newcommand{\capcalera}[4]{%
  \gdef\fm@cursnom{#2}%
  \par\noindent
  \begin{tcolorbox}[colback=curs,colframe=curs,boxrule=0pt,arc=2.2mm,width=3.5cm,height=2.15cm,
      valign=center,left=3mm,right=2mm,top=1mm,bottom=1mm,nobeforeafter,box align=center]
    \color{curstext}{\footnotesize #1}\par\vspace{1pt}{\fontsize{19}{21}\selectfont\bfseries #2\par}%
  \end{tcolorbox}\hspace{5mm}%
  \begin{minipage}[c]{\dimexpr\linewidth-3.5cm-5mm-2.9cm\relax}
    {\color{cursfosc}\small\bfseries\MakeUppercase{#3}}\par\vspace{2pt}
    {\color{navy}\fontsize{20}{23}\selectfont\bfseries #4\par}
  \end{minipage}\hfill
  \begin{minipage}[c]{2.6cm}\raggedleft\includegraphics[width=2.6cm]{\fm@logo}\end{minipage}\par
  \vspace{2pt}}

% ---------- barres de secció del color del curs ----------
\renewcommand{\seccio}[2]{\par\needspace{\fmneed\baselineskip}\bigskip
  \begin{tcolorbox}[colback=curs,colframe=curs,coltext=curstext,arc=2mm,boxrule=0pt,left=3mm,right=3mm,top=1.2mm,bottom=1.2mm]
  {\large\bfseries #1\quad #2}
  \end{tcolorbox}\vspace{-2pt}}

% ---------- estrelles de dificultat del color del curs (a Batxillerat no se'n posen) ----------
\@ifundefined{estrelles}{}{%
  \renewcommand{\estrelles}[1]{{\color{curs}\ifcase#1\or\faStar\faStar[regular]\faStar[regular]\or\faStar\faStar\faStar[regular]\or\faStar\faStar\faStar\fi}}}

% ---------- capçalera i peu de les pàgines ----------
% primera pàgina: logos institucionals al peu
\fancypagestyle{primera}{\fancyhf{}\renewcommand{\headrulewidth}{0pt}
  \fancyfoot[L]{\raisebox{-0.32\height}{\includegraphics[height=0.78cm]{logo-gva.png}}\hspace{5mm}%
    \raisebox{-0.32\height}{\includegraphics[height=0.78cm]{logo-vermellar.png}}}
  \fancyfoot[R]{\scriptsize\color{muted}\href{https://creativecommons.org/licenses/by-nc-sa/4.0/deed.ca}{\ccbyncsa}\ \ CC BY-NC-SA 4.0 · fentmates.org\qquad\small\thepage}}
% la resta: el símbol a l'esquerra i el curs a la dreta, amb el títol al centre com abans
\fancyhead[L]{\raisebox{-0.3\height}{\includegraphics[height=0.62cm]{\fm@marca}}}
\fancyhead[R]{\small\bfseries\color{cursfosc}\fm@cursnom}
\renewcommand{\headrule}{\color{curs!45}\hrule height \headrulewidth}
\makeatother

%   \fmcurs{1eso}
% i, al començament del document:
%   \capcalera{Matemàtiques}{1r ESO}{Fitxa d'exercicis · Unitat 2}{Divisibilitat}
% Opcional abans de \input: \def\fmneed{14} (espai mínim abans de cada secció).
% =====================================================================
\makeatletter
\def\fm@logodir{../../logo/pdf}
\providecommand{\fmneed}{10}
\def\fm@cursnom{}

% clau del curs, color i color del text damunt del color
\newcommand{\fm@curs}[3]{\@namedef{fm@color@#1}{#2}\@namedef{fm@text@#1}{#3}}
\fm@curs{1eso}{E6B84A}{navy}
\fm@curs{2eso}{3B82F6}{white}
\fm@curs{3eso}{FF7A50}{white}
\fm@curs{4eso-a}{A78BFA}{white}
\fm@curs{4eso-b}{6D28D9}{white}
\fm@curs{1batx-mi}{9B1C1C}{white}
\fm@curs{1batx-ccss}{EF6B6B}{white}
\fm@curs{2batx-mii}{1E6B47}{white}
\fm@curs{2batx-ccss}{2BB3A3}{white}

\newcommand{\fmcurs}[1]{%
  \@ifundefined{fm@color@#1}{\PackageError{fentmates}{Curs desconegut: #1}{}}{}%
  \edef\fm@tmp{\noexpand\definecolor{curs}{HTML}{\@nameuse{fm@color@#1}}%
               \noexpand\colorlet{curstext}{\@nameuse{fm@text@#1}}}\fm@tmp
  \colorlet{cursfosc}{curs!72!black}%
  \def\fm@logo{\fm@logodir/vertical-#1.pdf}%
  \def\fm@marca{\fm@logodir/marca-#1.pdf}}

% ---------- capçalera de la primera pàgina ----------
% #1 matèria (petit), #2 curs, #3 tipus de document i unitat, #4 títol
\newcommand{\capcalera}[4]{%
  \gdef\fm@cursnom{#2}%
  \par\noindent
  \begin{tcolorbox}[colback=curs,colframe=curs,boxrule=0pt,arc=2.2mm,width=3.5cm,height=2.15cm,
      valign=center,left=3mm,right=2mm,top=1mm,bottom=1mm,nobeforeafter,box align=center]
    \color{curstext}{\footnotesize #1}\par\vspace{1pt}{\fontsize{19}{21}\selectfont\bfseries #2\par}%
  \end{tcolorbox}\hspace{5mm}%
  \begin{minipage}[c]{\dimexpr\linewidth-3.5cm-5mm-2.9cm\relax}
    {\color{cursfosc}\small\bfseries\MakeUppercase{#3}}\par\vspace{2pt}
    {\color{navy}\fontsize{20}{23}\selectfont\bfseries #4\par}
  \end{minipage}\hfill
  \begin{minipage}[c]{2.6cm}\raggedleft\includegraphics[width=2.6cm]{\fm@logo}\end{minipage}\par
  \vspace{2pt}}

% ---------- barres de secció del color del curs ----------
\renewcommand{\seccio}[2]{\par\needspace{\fmneed\baselineskip}\bigskip
  \begin{tcolorbox}[colback=curs,colframe=curs,coltext=curstext,arc=2mm,boxrule=0pt,left=3mm,right=3mm,top=1.2mm,bottom=1.2mm]
  {\large\bfseries #1\quad #2}
  \end{tcolorbox}\vspace{-2pt}}

% ---------- estrelles de dificultat del color del curs (a Batxillerat no se'n posen) ----------
\@ifundefined{estrelles}{}{%
  \renewcommand{\estrelles}[1]{{\color{curs}\ifcase#1\or\faStar\faStar[regular]\faStar[regular]\or\faStar\faStar\faStar[regular]\or\faStar\faStar\faStar\fi}}}

% ---------- capçalera i peu de les pàgines ----------
% primera pàgina: logos institucionals al peu
\fancypagestyle{primera}{\fancyhf{}\renewcommand{\headrulewidth}{0pt}
  \fancyfoot[L]{\raisebox{-0.32\height}{\includegraphics[height=0.78cm]{logo-gva.png}}\hspace{5mm}%
    \raisebox{-0.32\height}{\includegraphics[height=0.78cm]{logo-vermellar.png}}}
  \fancyfoot[R]{\scriptsize\color{muted}\href{https://creativecommons.org/licenses/by-nc-sa/4.0/deed.ca}{\ccbyncsa}\ \ CC BY-NC-SA 4.0 · fentmates.org\qquad\small\thepage}}
% la resta: el símbol a l'esquerra i el curs a la dreta, amb el títol al centre com abans
\fancyhead[L]{\raisebox{-0.3\height}{\includegraphics[height=0.62cm]{\fm@marca}}}
\fancyhead[R]{\small\bfseries\color{cursfosc}\fm@cursnom}
\renewcommand{\headrule}{\color{curs!45}\hrule height \headrulewidth}
\makeatother

\fmcurs{3eso}

\begin{document}
\thispagestyle{primera}

% ---------- encabezado ----------
\capcalera{Matemáticas}{3\textsuperscript{o} ESO}{Ficha de ejercicios · Unidad 3}{Proporcionalidad y porcentajes}
\vspace{10pt}
Nombre y apellidos: \rule{6.6cm}{0.4pt}\hfill Grupo: \rule{1.6cm}{0.4pt}\hfill Fecha: \rule{2.2cm}{0.4pt}
\vspace{8pt}

\begin{tcolorbox}[colback=soft,colframe=soft,arc=2mm,left=3mm,right=3mm,top=1.5mm,bottom=1.5mm]
\small Ejercicios graduados por niveles de dificultad:\quad
\estrelles{1}\ básico \qquad \estrelles{2}\ medio \qquad \estrelles{3}\ avanzado.\\[2pt]
Resuélvelos en el cuaderno mostrando todos los pasos. Las tablas se pueden completar directamente en la ficha. Puedes usar la calculadora.
\end{tcolorbox}


% =====================================================
\seccio{1}{Porcentajes, fracciones y decimales}

\ex{1} Completa la tabla:\par\smallskip
{\renewcommand{\arraystretch}{1.75}\centering
\begin{tabular}{|C{4.2cm}|C{4.2cm}|C{4.2cm}|}
\hline
\rowcolor{soft}\textbf{Porcentaje} & \textbf{Fracción irreducible} & \textbf{Decimal} \\ \hline
$40\,\%$ & & \\ \hline
 & $\dfrac34$ & \\ \hline
 & & $0,15$ \\ \hline
$8\,\%$ & & \\ \hline
 & $\dfrac{7}{20}$ & \\ \hline
 & & $1,25$ \\ \hline
$2,5\,\%$ & & \\ \hline
\end{tabular}\par}

\ex{1} Expresa como porcentaje:
\begin{tasks}(6)
  \task $0,32$
  \task $0,07$
  \task $1,5$
  \task $\dfrac35$
  \task $\dfrac{9}{25}$
  \task $\dfrac18$
\end{tasks}

\ex{2} Escribe como fracción irreducible:
\begin{tasks}(6)
  \task $12\,\%$
  \task $65\,\%$
  \task $0,8\,\%$
  \task $12,5\,\%$
  \task $140\,\%$
  \task $0,25\,\%$
\end{tasks}

\ex{2} Ordena de menor a mayor:
\[
\frac25\qquad 0,38\qquad 41\,\%\qquad \frac{7}{16}\qquad 0,\overline{4}
\]

\ex{3} En el grupo A han aprobado 18 de los 24 alumnos; en el B, 21 de 30, y en el C, 16 de 20. ¿Qué porcentaje de aprobados tiene cada grupo? Ordena los grupos de mejor a peor resultado.

% =====================================================
\seccio{2}{Cálculo de porcentajes}

\ex{1} Calcula:
\begin{tasks}(3)
  \task El $20\,\%$ de $150$
  \task El $75\,\%$ de $64$
  \task El $5\,\%$ de $360$
  \task El $12\,\%$ de $250$
  \task El $150\,\%$ de $40$
  \task El $0,5\,\%$ de $2\,000$
\end{tasks}

\ex{1} Calcula mentalmente:
\begin{tasks}(3)
  \task El $10\,\%$ de $470$
  \task El $50\,\%$ de $86$
  \task El $25\,\%$ de $120$
  \task El $1\,\%$ de $3\,500$
  \task El $20\,\%$ de $45$
  \task El $15\,\%$ de $80$
\end{tasks}

\ex{2} En un instituto de $850$ alumnos, el $46\,\%$ son chicas. ¿Cuántas chicas y cuántos chicos hay?

\ex{2} Un jersey de $45$ € tiene un descuento del $30\,\%$. ¿Cuántos euros nos descuentan? ¿Cuánto pagamos?

\ex{2} Se ha preguntado a $1\,200$ personas por su deporte preferido: el $35\,\%$ ha elegido el fútbol, el $28\,\%$ el baloncesto, el $15\,\%$ el tenis y el resto, otros deportes. ¿Cuántas personas han elegido cada opción?

\ex{3} Una empresa tiene un presupuesto de $2\,500$ €. Gasta el $40\,\%$ en ordenadores y el $30\,\%$ de lo que le queda en programas. El resto se dedica a formación. ¿Cuántos euros se destinan a cada cosa? ¿Qué porcentaje del presupuesto total es la formación?

\ex{3} Calcula el $20\,\%$ del $30\,\%$ de $500$. ¿Es lo mismo que el $30\,\%$ del $20\,\%$ de $500$? ¿A qué porcentaje único de $500$ equivale?

% =====================================================
\seccio{3}{Tanto por ciento que representa un número respecto a otro}

\ex{1} ¿Qué porcentaje representa:
\begin{tasks}(3)
  \task $12$ respecto a $48$?
  \task $36$ respecto a $120$?
  \task $9$ respecto a $12$?
  \task $45$ respecto a $30$?
  \task $3$ respecto a $200$?
  \task $7$ respecto a $35$?
\end{tasks}

\ex{1} En una clase de $25$ alumnos han aprobado $21$. ¿Qué porcentaje ha aprobado? ¿Y cuál ha suspendido?

\ex{2} Calcula el número en cada caso:
\begin{tasks}(2)
  \task El $15\,\%$ de un número es $27$.
  \task El $40\,\%$ de un número es $18$.
  \task El $120\,\%$ de un número es $54$.
  \task El $2,5\,\%$ de un número es $5$.
\end{tasks}

\ex{2} Un jugador ha marcado $18$ de los $24$ penaltis que ha lanzado y otro, $22$ de $30$. ¿Quién tiene mejor porcentaje de acierto?

\ex{2} Para reservar un viaje pagamos una entrada de $36$ €, que es el $12\,\%$ del precio total. ¿Cuánto cuesta el viaje? ¿Cuánto nos queda por pagar?

\ex{3} Un embalse tiene una capacidad de $450$ hm$^3$ y ahora contiene $162$ hm$^3$.
\begin{tasks}(1)
  \task ¿A qué porcentaje de su capacidad está?
  \task Después de unas lluvias llega al $40\,\%$. ¿Cuántos hm$^3$ de agua ha ganado?
\end{tasks}

\ex{3} En una ciudad, $\frac38$ de la población son menores de $30$ años y, de estos, el $40\,\%$ estudia. ¿Qué porcentaje de la población total son estudiantes menores de $30$ años?

% =====================================================
\seccio{4}{Aumentos y disminuciones porcentuales}

\ex{1} Completa la tabla:\par\smallskip
{\renewcommand{\arraystretch}{1.6}\centering
\begin{tabular}{|C{3.2cm}|C{3.2cm}|C{3.2cm}|C{3.2cm}|}
\hline
\rowcolor{soft}\textbf{Cantidad inicial} & \textbf{Variación} & \textbf{Índice de variación} & \textbf{Cantidad final} \\ \hline
$200$ & aumento del $15\,\%$ & & \\ \hline
$80$ & disminución del $25\,\%$ & & \\ \hline
$1\,200$ & aumento del $3\,\%$ & & \\ \hline
$500$ & & $0,94$ & \\ \hline
$64$ & & & $80$ \\ \hline
\end{tabular}\par}

\ex{1} ¿Qué aumento o disminución porcentual indica cada índice de variación?
\begin{tasks}(6)
  \task $1,12$
  \task $0,85$
  \task $1,03$
  \task $0,6$
  \task $2,5$
  \task $0,995$
\end{tasks}

\ex{2} Un ordenador cuesta $640$ € sin IVA. ¿Cuál es el precio final si el IVA es del $21\,\%$?

\ex{2} El alquiler de un piso es de $750$ € al mes y sube un $3,2\,\%$. ¿Cuál es el nuevo alquiler?

\ex{2} La población de un pueblo ha pasado de $12\,500$ a $11\,750$ habitantes. ¿Qué porcentaje ha variado?

\ex{3} Una tienda sube todos los precios un $10\,\%$ y otra los rebaja todos un $10\,\%$. Una misma chaqueta costaba $90$ € en la primera tienda y $110$ € en la segunda. ¿En cuál es más barata ahora?

\ex{3} Un abono al gimnasio cuesta $45$ € al mes. Si se paga el año entero de una vez, se hace un descuento del $15\,\%$ sobre el precio de los $12$ meses. ¿Cuánto se paga con el abono anual? ¿Cuánto se ahorra?

% =====================================================
\seccio{5}{Cálculo de la cantidad inicial}

\ex{1} Calcula la cantidad inicial:
\begin{tasks}(2)
  \task Después de aumentar un $20\,\%$ vale $84$.
  \task Después de disminuir un $30\,\%$ vale $56$.
  \task Después de aumentar un $5\,\%$ vale $420$.
  \task Después de disminuir un $12\,\%$ vale $220$.
\end{tasks}

\ex{1} Con el $21\,\%$ de IVA incluido, unos auriculares cuestan $181,50$ €. ¿Cuál es el precio sin IVA?

\ex{2} Rebajada un $35\,\%$, una bicicleta cuesta $520$ €. ¿Cuál era el precio original? ¿Cuánto nos ahorramos?

\ex{2} Después de una subida del $3\,\%$, una trabajadora cobra $1\,802,50$ €. ¿Cuál era su sueldo antes de la subida?

\ex{2} Un coche de ocasión se vende por $9\,600$ €, que es un $40\,\%$ menos de lo que valía nuevo. ¿Cuánto valía nuevo?

\ex{3} Una factura de $96,80$ € incluye el $21\,\%$ de IVA. Luis dice: «El IVA que hemos pagado es el $21\,\%$ de $96,80$, es decir, $20,33$ €».
\begin{tasks}(1)
  \task Calcula el importe sin IVA y el IVA que se ha pagado realmente.
  \task Explica cuál es el error de Luis.
\end{tasks}

\ex{3} Un depósito de $4\,000$ L de capacidad ha perdido el $8\,\%$ del agua que tenía y ahora contiene $3\,450$ L. ¿Cuántos litros tenía antes de la pérdida? ¿A qué porcentaje de su capacidad estaba?

% =====================================================
\seccio{6}{Variaciones porcentuales encadenadas}

\ex{1} Calcula el índice de variación total y el porcentaje de variación global:
\begin{tasks}(2)
  \task Un aumento del $10\,\%$ y después uno del $20\,\%$.
  \task Una disminución del $10\,\%$ y después una del $20\,\%$.
  \task Un aumento del $50\,\%$ y después una disminución del $50\,\%$.
  \task Un aumento del $25\,\%$ y después una disminución del $20\,\%$.
\end{tasks}

\ex{1} Un precio de $300$ € sube un $5\,\%$ y, más tarde, un $8\,\%$. ¿Cuál es el precio final?

\ex{2} Un artículo de $120$ € está rebajado un $30\,\%$ y en la caja nos hacen un $10\,\%$ de descuento adicional. ¿Cuánto pagamos? ¿Cuál es el descuento total que nos han hecho?

\ex{2} El precio de la gasolina sube un $4\,\%$, después baja un $6\,\%$ y, finalmente, sube un $2\,\%$. ¿Ha subido o ha bajado en total? ¿Qué porcentaje?

\ex{2} Invertimos $5\,000$ € en bolsa. El primer año ganan un $12\,\%$ y el segundo pierden un $15\,\%$. ¿Cuánto dinero tenemos al final? ¿Hemos ganado o perdido? ¿Qué porcentaje?

\ex{3} Un precio sube un $30\,\%$. ¿Qué porcentaje tiene que bajar para volver al precio inicial?

\ex{3} Un precio sube un $20\,\%$ y después baja un cierto porcentaje. El resultado es el mismo que si hubiera bajado directamente un $4\,\%$. ¿Qué porcentaje ha bajado la segunda vez?

% =====================================================
\seccio{7}{Interés compuesto}

\ex{1} Depositamos $3\,000$ € al $4\,\%$ de interés compuesto anual durante $2$ años. ¿Qué capital tendremos? ¿Cuántos intereses habremos ganado?

\ex{1} Completa la tabla de un capital de $1\,000$ € al $5\,\%$ de interés compuesto anual:\par\smallskip
{\renewcommand{\arraystretch}{1.6}\centering
\begin{tabular}{|C{2.2cm}|C{4cm}|C{4cm}|C{4cm}|}
\hline
\rowcolor{soft}\textbf{Año} & \textbf{Capital al inicio} & \textbf{Intereses del año} & \textbf{Capital al final} \\ \hline
$1$ & $1\,000$ & & \\ \hline
$2$ & & & \\ \hline
$3$ & & & \\ \hline
\end{tabular}\par}

\ex{2} Calcula el capital final que se obtiene al invertir $8\,000$ € al $2,5\,\%$ de interés compuesto durante $5$ años. ¿Cuántos intereses se obtendrían con interés simple?

\ex{2} ¿Qué capital tenemos que depositar al $3\,\%$ de interés compuesto anual para tener $10\,927,27$ € dentro de $3$ años?

\ex{2} Depositamos $12\,000$ € al $3\,\%$ anual con abono mensual de los intereses durante $2$ años. ¿Qué capital tendremos?

\ex{3} ¿Cuántos años tienen que pasar para que $5\,000$ € al $4\,\%$ de interés compuesto anual superen los $6\,000$ €?

\ex{3} Queremos invertir $10\,000$ € durante $4$ años. El banco A nos ofrece un $3,5\,\%$ anual con abono anual y el banco B, un $3,4\,\%$ anual con abono mensual. ¿Qué oferta es mejor? ¿Qué diferencia hay?

% =====================================================
\seccio{8}{Repartos proporcionales}

\ex{1} Reparte $720$ de manera directamente proporcional a $3$, $4$ y $5$.

\ex{1} Tres hermanos han trabajado $10$, $15$ y $20$ horas en un negocio familiar y se reparten $1\,800$ € proporcionalmente a las horas trabajadas. ¿Cuánto recibe cada uno?

\ex{2} Reparte $1\,080$ de manera inversamente proporcional a $3$, $4$ y $6$.

\ex{2} Un abuelo deja $9\,000$ € a sus tres nietos para que se los repartan de manera inversamente proporcional a sus edades: $6$, $9$ y $18$ años. ¿Cuánto recibe cada uno?

\ex{2} Tres socios han invertido $15\,000$ €, $20\,000$ € y $25\,000$ € en una empresa. Si los beneficios han sido de $9\,600$ €, ¿cuánto le corresponde a cada uno?

\ex{3} Tres socios invierten en un negocio: A pone $6\,000$ € durante $8$ meses; B, $4\,000$ € durante $12$ meses, y C, $10\,000$ € durante $3$ meses. ¿Cómo se tienen que repartir unos beneficios de $8\,400$ €?

\ex{3} Una prima de $4\,500$ € se reparte entre dos trabajadores de manera directamente proporcional a las horas extra que han hecho ($30$ y $40$) e inversamente proporcional a las faltas que han tenido ($3$ y $2$). ¿Cuánto recibe cada uno?

\ex{3} Un premio se reparte de manera directamente proporcional a $2$, $3$ y $5$. Si el tercero recibe $450$ € más que el primero, ¿de cuánto es el premio? ¿Cuánto recibe cada uno?

% =====================================================
\seccio{9}{Problemas de mezclas}

\ex{1} Se mezclan $5$ kg de café de $6$ €/kg con $3$ kg de café de $10$ €/kg. ¿A qué precio sale el kilo de mezcla?

\ex{1} Se mezclan $12$ L de vino de $2,50$ €/L con $8$ L de vino de $4$ €/L. ¿Cuál es el precio del litro de mezcla?

\ex{2} ¿Cuántos kilos de café de $9$ €/kg se tienen que mezclar con $20$ kg de café de $6$ €/kg para obtener una mezcla de $7$ €/kg?

\ex{2} Se mezclan $3$ L de un zumo con un $60\,\%$ de fruta y $2$ L de una bebida con un $20\,\%$ de fruta. ¿Qué porcentaje de fruta tiene la mezcla?

\ex{2} Un comerciante mezcla $30$ kg de arroz de $1,80$ €/kg con $20$ kg de arroz de $2,40$ €/kg. ¿A qué precio tiene que vender el kilo de mezcla para ganar un $25\,\%$?

\ex{3} Se funden $40$ g de oro de ley $0,750$ con $10$ g de oro puro. ¿Cuál es la ley de la aleación obtenida? ¿Cuántos quilates tiene?

\ex{3} ¿Cuántos gramos de oro puro se tienen que añadir a $60$ g de oro de $18$ quilates para obtener oro de $20$ quilates?

\ex{3} Tenemos $15$ L de alcohol al $70\,\%$. ¿Cuánta agua hay que añadir para que quede al $40\,\%$?

% =====================================================
\seccio{10}{Problemas de movimientos}

\ex{1} Expresa $54$ km/h en m/s y $25$ m/s en km/h.

\ex{1} Un tren circula a $120$ km/h.
\begin{tasks}(1)
  \task ¿Cuánto tiempo tarda en recorrer $330$ km? Exprésalo en horas y minutos.
  \task ¿Qué distancia recorre en $1$ h $20$ min?
\end{tasks}

\ex{2} Dos ciudades están a $360$ km. A las $9{:}00$ sale un coche de cada ciudad, el uno hacia el otro, a $70$ km/h y a $90$ km/h. ¿A qué hora se encuentran? ¿A qué distancia de la ciudad del primer coche?

\ex{2} Un peatón sale a $5$ km/h. Dos horas más tarde, sale del mismo lugar y por el mismo camino un ciclista a $15$ km/h. ¿Cuánto tiempo tarda en alcanzarlo? ¿A qué distancia del punto de partida?

\ex{2} Dos ciclistas salen a la vez del mismo punto en sentidos contrarios, a $18$ km/h y a $22$ km/h. ¿Cuánto tiempo pasará hasta que estén a $100$ km el uno del otro?

\ex{3} Las ciudades A y B están a $500$ km. A las $7{:}30$ sale un camión de A hacia B a $80$ km/h y a las $9{:}00$ sale un coche de B hacia A a $110$ km/h. ¿A qué hora se cruzan? ¿A qué distancia de A?

\ex{3} Un tren de mercancías sale de una estación a $60$ km/h. Una hora y media más tarde, sale de la misma estación, por una vía paralela y en el mismo sentido, un tren rápido a $150$ km/h. ¿Cuánto tarda el rápido en alcanzar al de mercancías? ¿A cuántos kilómetros de la estación?

\ex{3} Dos atletas salen a la vez y en el mismo sentido en una pista de $400$ m, a $6$ m/s y a $5$ m/s. ¿Cuánto tiempo tarda la más rápida en doblar a la otra? ¿Cuántas vueltas ha dado cada una en ese tiempo?

% =====================================================
\seccio{11}{Proporcionalidad compuesta}

\ex{1} Indica si cada pareja de magnitudes es directamente proporcional, inversamente proporcional o no proporcional:
\begin{tasks}(2)
  \task La velocidad de un coche y el tiempo que tarda en hacer un trayecto.
  \task Los kilos de fruta que compramos y el precio que pagamos.
  \task La edad de una persona y su altura.
  \task El número de obreros y los días que tardan en hacer una obra.
  \task El lado de un cuadrado y su área.
  \task El número de grifos abiertos y los litros que vierten en una hora.
\end{tasks}

\ex{1} Resuelve:
\begin{tasks}(1)
  \task Si $4$ entradas de cine cuestan $26$ €, ¿cuánto cuestan $7$ entradas?
  \task $8$ operarios hacen un trabajo en $15$ días. ¿Cuántos días tardarían $12$ operarios?
\end{tasks}

\ex{2} $6$ grifos abiertos durante $4$ horas vierten $3\,600$ L de agua. ¿Cuántos litros verterán $9$ grifos abiertos durante $5$ horas?

\ex{2} $10$ vacas comen $1\,800$ kg de forraje en $12$ días. ¿Cuánto forraje comerán $15$ vacas en $20$ días?

\ex{2} $5$ obreros, trabajando $8$ horas diarias, hacen una obra en $18$ días. ¿Cuántos días tardarían $6$ obreros trabajando $10$ horas diarias?

\ex{3} $12$ impresoras, funcionando $6$ horas diarias, imprimen $54\,000$ hojas en $5$ días. ¿Cuántos días necesitarán $10$ impresoras, funcionando $8$ horas diarias, para imprimir $72\,000$ hojas?

\ex{3} Con $120$ kg de pienso se alimentan $30$ gallinas durante $20$ días. ¿Cuántas gallinas se pueden alimentar durante $25$ días con $90$ kg de pienso?

\ex{3} $8$ obreros, trabajando $9$ horas diarias, construyen $120$ m de muro en $15$ días. ¿Cuántos días tardarán $10$ obreros, trabajando $6$ horas diarias, en construir $200$ m de muro?

\vfill
\begin{center}\color{muted}\small\faCheckCircle[regular]\ \ ¡Ánimo! Si te cuesta algún ejercicio, repasa los apuntes interactivos de la unidad.\end{center}

\end{document}
